\documentclass{article}
\usepackage{amsmath,amssymb}
\usepackage{graphicx}
\usepackage{hyperref}
\usepackage{geometry}
\geometry{margin=1in}

\begin{document}

\title{
Chain Explosion Model:\\
A Local Discrete Propagation Model for Interference, Measurement, and Entanglement
}

\author{
Independent Researcher \\
\href{https://github.com/tomnattle/chain-explosion-model}{github.com/tomnattle/chain-explosion-model}
}

\date{\today}

\maketitle

\begin{abstract}
We introduce a local, discrete, causal lattice propagation model termed the \emph{chain explosion model}.
Energy propagates on a 2D grid using nearest-neighbor and diagonal coupling with directional strength
parameters and uniform attenuation. The model naturally reproduces double-slit interference,
which-way measurement via local absorption, delayed-choice-like behavior, and entangled wave-packet
correlations without nonlocality or wavefunction collapse.

A unique, testable prediction of the model is that \emph{interference visibility decays with propagation distance},
while phase correlations between split wave packets remain stable even as amplitude decays exponentially.
This separates the model from standard quantum mechanics and classical wave optics.
All simulations are reproducible via open-source Python code.
\end{abstract}

\section{Introduction}
Standard quantum mechanics exhibits highly nonintuitive features including nonlocality,
wavefunction collapse, and entanglement that appears independent of distance.
Many attempts have been made to construct underlying realistic models,
including pilot-wave theories and lattice-based cellular automata.

This work presents a purely discrete, local, causal model:
the \emph{chain explosion model}. It is defined entirely by
grid-based energy propagation rules, without Hilbert spaces,
complex amplitudes, or nonlocal interactions.

Despite its simplicity, the model reproduces key phenomena:
\begin{itemize}
\item Double-slit interference fringes
\item Visibility reduction due to which-way measurement
\item Delayed-choice behavior
\item Entangled-like correlations from symmetric wave-packet splitting
\item Stable phase correlation with decaying amplitude
\end{itemize}

\section{Model Definition}
The model evolves on a 2D lattice $(x,y)$ with a real-valued energy field $E(x,y)$.
At each time step, energy propagates to neighboring lattice sites
with fixed directional coupling strengths:
\begin{align*}
A &: \text{forward (right) coupling} \\
S &: \text{lateral (up/down) coupling} \\
B &: \text{backward (left) coupling}
\end{align*}
An attenuation factor $\lambda \in (0,1)$ is applied at each step.

Propagation rule summary:
\begin{enumerate}
\item Energy at each site is attenuated by $\lambda$.
\item Energy is distributed to adjacent and diagonal sites according to $A,S,B$.
\item Barriers block propagation; slits allow transmission.
\item Detectors/measurements are implemented as local absorption.
\item Split wave packets maintain consistent phase evolution without decay.
\end{enumerate}

The model is \emph{strictly local}:
no site influences any other site except through step-by-step propagation.

\section{Double-Slit Interference}
A barrier with two slits is placed at fixed $x$.
The initial source emits a narrow beam.
Interference fringes form naturally on a distant screen.

The model predicts a novel effect:
\textbf{interference visibility decreases with increasing screen distance}.
This is not predicted by standard quantum or classical optics
and provides a direct experimental test.

Visibility is defined as:
\[
V = \frac{I_{\text{max}} - I_{\text{min}}}{I_{\text{max}} + I_{\text{min}}}
\]

Simulations show $V$ decreases smoothly as distance increases,
while total energy on the screen decays exponentially.

\section{Measurement and Which-Wat Information}
Measurement is modeled as \emph{local absorption} at one slit.
Three regimes are simulated:
\begin{itemize}
\item No absorption: full interference
\item Partial absorption: reduced visibility
\item Full absorption: no interference
\end{itemize}

No collapse postulate is needed.
Decoherence emerges purely from local energy removal.

\section{Delayed-Choice Behavior}
Absorption is activated \emph{after} the wave front has passed the slits.
The model still shows suppression of interference,
consistent with delayed-choice experiments.

This occurs without retrocausality:
absorption removes energy from the propagating front,
erasing the superposition required for fringes.

\section{Entanglement-Like Correlations}
A symmetric beam splitter divides the wave into two spatially separated packets.
Two distant detectors measure energy and phase.

Key results:
\begin{itemize}
\item Energy at detectors decays exponentially with distance
\item Phase correlation remains nearly perfect ($\cos(\Delta\phi)\approx 1$)
\item No nonlocal interaction is required
\end{itemize}

The model explains entanglement as \emph{shared origin and consistent phase evolution},
not as nonlocal connection.

\section{Unique Testable Predictions}
The chain explosion model makes predictions that distinguish it from standard theory:
\begin{enumerate}
\item Interference visibility decreases with propagation distance
\item Entangled signal amplitude decays, but phase correlation does not
\item Measurement is local absorption; no collapse needed
\end{enumerate}

These can be tested in precision optical experiments.

\section{Code and Reproducibility}
All simulations are open-source Python code available at:
\\
\href{https://github.com/tomnattle/chain-explosion-model}{https://github.com/tomnattle/chain-explosion-model}

Dependencies: numpy, matplotlib, numba (for accelerated scans).

\section{Conclusion}
The chain explosion model provides a local, realistic, discrete framework
that reproduces interference, measurement, delayed choice, and entanglement-like correlations.
It eliminates nonlocality and wavefunction collapse while making novel, testable predictions.

The model suggests that quantum phenomena may emerge from
simple local propagation rules at a discrete underlying level.

\end{document}